
\subsection{Übergang zur Quantenfeldtheorie}
\label{subsec:uebergang_zur_quantenfeldtheorie}

Die Dirac-Gleichung\index{Dirac-Gleichung!Grenzen der Einteilchentheorie}
beschreibt ein relativistisches Elektron mit Spin und führte zur
Vorhersage des Positrons. Damit löste sie wesentliche Probleme der
nichtrelativistischen Quantenmechanik.

Gleichzeitig zeigte sich jedoch eine neue Grenze: In relativistischen
Prozessen kann die Anzahl der Teilchen nicht immer als fest
vorausgesetzt werden.

Elektronen und Positronen können erzeugt und vernichtet werden. Eine
Theorie, die lediglich die Wellenfunktion eines einzelnen Elektrons
beschreibt, reicht deshalb für solche Vorgänge nicht mehr aus.

Der notwendige nächste Schritt führt zur
Quantenfeldtheorie\index{Quantenfeldtheorie}.

\begin{DidacticBox}[Die entscheidende Grenze der Einteilchentheorie]
	\label{box:grenze_einteilchentheorie}
	
	Die Dirac-Gleichung kann die Bewegung und den Spin eines einzelnen
	Elektrons beschreiben.
	
	Bei genügend hohen Energien können jedoch neue Elektronen und
	Positronen entstehen. Ebenso können Teilchen durch Annihilation
	verschwinden.
	
	Sobald sich die Anzahl der Teilchen ändern kann, genügt eine
	gewöhnliche Einteilchen-Wellenfunktion nicht mehr.
\end{DidacticBox}

\subsubsection*{Warum eine feste Teilchenzahl nicht ausreicht}

In der gewöhnlichen Quantenmechanik wird häufig ein System mit einer
festen Anzahl von Teilchen betrachtet.

Ein Elektron kann beispielsweise

\begin{itemize}
	\item seinen Ort verändern,
	\item seinen Impuls verändern,
	\item seinen Energiezustand wechseln oder
	\item mit einem äußeren Feld wechselwirken.
\end{itemize}

Es bleibt dabei jedoch als Elektron im betrachteten System vorhanden.

In der relativistischen Physik kann genügend Energie in neue Teilchen
umgewandelt werden. Ein Photon kann in der Nähe eines Atomkerns ein
Elektron-Positron-Paar erzeugen:

\begin{equation}
	\gamma
	+
	Z
	\longrightarrow
	e^-
	+
	e^+
	+
	Z.
	\label{eq:paarbildung_uebergang_feldtheorie}
\end{equation}

Umgekehrt können Elektron und Positron annihilieren:

\begin{equation}
	e^-
	+
	e^+
	\longrightarrow
	\gamma
	+
	\gamma.
	\label{eq:annihilation_uebergang_feldtheorie}
\end{equation}

Vor und nach diesen Prozessen ist die Anzahl der Teilchen
unterschiedlich.

Eine Theorie mit einer fest vorgegebenen Teilchenzahl kann solche
Vorgänge nicht vollständig erfassen.

\begin{PhysicsBox}[Veränderliche Teilchenzahl]
	\label{box:veraenderliche_teilchenzahl}
	
	In einer relativistischen Quantentheorie können
	
	\begin{itemize}
		\item Teilchen erzeugt,
		\item Antiteilchen erzeugt,
		\item Teilchen vernichtet und
		\item Strahlung in Materie oder Materie in Strahlung umgewandelt
		werden.
	\end{itemize}
	
	Die Erhaltung der Energie bleibt dabei bestehen. Energie kann jedoch
	in Ruhemasse, Bewegungsenergie und Strahlungsenergie übergehen.
\end{PhysicsBox}

\subsubsection*{Das Feld als grundlegende Größe}

Die Quantenfeldtheorie verändert den Ausgangspunkt der Beschreibung.

Nicht mehr das einzelne Teilchen steht im Mittelpunkt, sondern ein
Feld\index{Feld!Quantenfeldtheorie}, das sich über Raum und Zeit
erstreckt.

Für jede Teilchenart wird ein entsprechendes Quantenfeld betrachtet.
Beim Elektron ist dies das Elektronenfeld%
\index{Elektronenfeld} oder genauer das Dirac-Feld%
\index{Dirac-Feld}.

Ein Elektron wird dann als Anregung dieses Feldes verstanden.

Auch ein Positron ist eine Anregung desselben Feldes, jedoch mit den
Eigenschaften des Antiteilchens.

\begin{DidacticBox}[Teilchen als Anregungen eines Feldes]
	\label{box:teilchen_als_feldanregungen}
	
	In der Quantenfeldtheorie ist ein Elektron kein völlig eigenständiges
	Objekt, das unabhängig von einem Feld existiert.
	
	Es ist ein bestimmter angeregter Zustand des Elektronenfeldes.
	
	Ein Positron ist eine andere mögliche Anregung desselben Feldes.
\end{DidacticBox}

Diese Sichtweise ähnelt in gewisser Weise einer Welle auf einer
Wasseroberfläche. Die Welle ist kein zusätzliches Material neben dem
Wasser, sondern ein angeregter Zustand der Wasseroberfläche.

Der Vergleich ist jedoch nur begrenzt gültig. Ein Quantenfeld ist kein
klassischer Stoff und benötigt kein materielles Trägermedium.

\subsubsection*{Vom Dirac-Spinor zum Feldoperator}

In der Einteilchentheorie ist

\begin{equation}
	\Psi(\mathbf{r},t)
\end{equation}

eine Wellenfunktion beziehungsweise ein Dirac-Spinor.

In der Quantenfeldtheorie wird daraus ein
Feldoperator\index{Feldoperator}

\begin{equation}
	\Psi(\mathbf{r},t)
	\longrightarrow
	\hat{\Psi}(\mathbf{r},t).
	\label{eq:dirac_spinor_zum_feldoperator}
\end{equation}

Der Operator \(\hat{\Psi}\) beschreibt nicht mehr nur den Zustand eines
bereits vorhandenen Elektrons. Er kann auch Zustände mit
unterschiedlichen Teilchenzahlen miteinander verbinden.

Vereinfacht enthält das quantisierte Dirac-Feld Anteile, die

\begin{itemize}
	\item ein Elektron vernichten und
	\item ein Positron erzeugen
\end{itemize}

können.

Der adjungierte Feldoperator enthält entsprechend Anteile, die

\begin{itemize}
	\item ein Elektron erzeugen und
	\item ein Positron vernichten
\end{itemize}

können.

\begin{MathBox}[Symbolische Struktur des Dirac-Feldes]
	\label{box:symbolische_struktur_dirac_feld}
	
	Vereinfacht kann die Struktur des quantisierten Dirac-Feldes als
	
	\begin{equation}
		\hat{\Psi}
		\sim
		\hat{b}
		+
		\hat{d}^{\dagger}
	\end{equation}
	
	geschrieben werden.
	
	Dabei steht
	
	\begin{itemize}
		\item \(\hat{b}\) für die Vernichtung eines Elektrons und
		\item \(\hat{d}^{\dagger}\) für die Erzeugung eines Positrons.
	\end{itemize}
	
	Die vollständige Gleichung enthält zusätzlich Impuls-, Spin- und
	Raumzeitabhängigkeiten.
\end{MathBox}

Die Symbole \(\hat{b}\) und \(\hat{d}^{\dagger}\) werden als
Vernichtungs- und Erzeugungsoperatoren%
\index{Vernichtungsoperator}%
\index{Erzeugungsoperator}
bezeichnet.

Sie verändern die Anzahl der Teilchen in einem Quantenzustand.

\subsubsection*{Zustände mit unterschiedlicher Teilchenzahl}

Der Zustandsraum einer Quantenfeldtheorie muss Zustände mit
unterschiedlich vielen Teilchen enthalten.

Ein möglicher Zustand kann beispielsweise kein Elektron enthalten:

\begin{equation}
	\lvert 0\rangle.
\end{equation}

Dieser Zustand wird als Vakuumzustand%
\index{Vakuumzustand} bezeichnet.

Ein Zustand mit einem Elektron kann symbolisch geschrieben werden als

\begin{equation}
	\lvert e^-\rangle.
\end{equation}

Ein Zustand mit einem Elektron und einem Positron lautet beispielsweise

\begin{equation}
	\lvert e^-e^+\rangle.
\end{equation}

Durch einen Erzeugungsoperator kann aus dem Vakuum ein Teilchenzustand
entstehen:

\begin{equation}
	\hat{b}^{\dagger}
	\lvert 0\rangle
	=
	\lvert e^-\rangle.
	\label{eq:erzeugung_elektron_vakuum}
\end{equation}

Entsprechend kann ein Vernichtungsoperator ein vorhandenes Elektron
entfernen:

\begin{equation}
	\hat{b}
	\lvert e^-\rangle
	=
	\lvert 0\rangle.
	\label{eq:vernichtung_elektron_vakuum}
\end{equation}

\begin{DidacticBox}[Was bedeutet Teilchenvernichtung?]
	\label{box:bedeutung_teilchenvernichtung}
	
	Teilchenvernichtung bedeutet nicht, dass Energie oder andere
	Erhaltungsgrößen verschwinden.
	
	Der ursprüngliche Teilchenzustand geht in einen anderen Zustand über.
	Die Energie kann beispielsweise in Photonen oder in die
	Bewegungsenergie anderer Teilchen übergehen.
	
	Die Gesamtenergie und die übrigen Erhaltungsgrößen bleiben erhalten.
\end{DidacticBox}

\subsubsection*{Das elektromagnetische Feld wird ebenfalls quantisiert}

Für eine vollständige Beschreibung der Wechselwirkung zwischen
Elektronen und Licht genügt es nicht, nur das Elektronenfeld zu
quantisieren.

Auch das elektromagnetische Feld%
\index{elektromagnetisches Feld!Quantisierung}
muss quantenmechanisch behandelt werden.

Die Anregungen des quantisierten elektromagnetischen Feldes sind
Photonen\index{Photon!Feldanregung}.

Damit ergibt sich eine einheitliche Beschreibung:

\begin{itemize}
	\item Elektronen sind Anregungen des Elektronenfeldes.
	\item Positronen sind Antiteilchenanregungen des Elektronenfeldes.
	\item Photonen sind Anregungen des elektromagnetischen Feldes.
\end{itemize}

Die Wechselwirkung zwischen diesen Feldern wird durch die
Quantenelektrodynamik%
\index{Quantenelektrodynamik} beschrieben.

\begin{PhysicsBox}[Grundidee der Quantenelektrodynamik]
	\label{box:grundidee_quantenelektrodynamik}
	
	Die Quantenelektrodynamik beschreibt die Wechselwirkung zwischen
	
	\begin{itemize}
		\item dem Elektronen-Positronen-Feld und
		\item dem elektromagnetischen Feld.
	\end{itemize}
	
	Dabei können Elektronen und Positronen Photonen aussenden oder
	aufnehmen.
	
	Die Quantenelektrodynamik wird häufig mit der Abkürzung QED bezeichnet.
\end{PhysicsBox}

\subsubsection*{Wechselwirkung zwischen Elektronen und Photonen}

Ein Elektron kann ein Photon aussenden:

\begin{equation}
	e^-
	\longrightarrow
	e^-
	+
	\gamma.
	\label{eq:elektron_emission_photon_symbolisch}
\end{equation}

Ein Elektron kann auch ein Photon absorbieren:

\begin{equation}
	e^-
	+
	\gamma
	\longrightarrow
	e^-.
	\label{eq:elektron_absorption_photon_symbolisch}
\end{equation}

Diese symbolischen Gleichungen zeigen lediglich die beteiligten
Teilchen. Ein freies Elektron kann im leeren Raum nicht beliebig ein
reales Photon aussenden oder aufnehmen, da dabei gleichzeitig Energie
und Impuls erhalten bleiben müssen.

In der Nähe eines Atomkerns, innerhalb eines Atoms oder bei der
Wechselwirkung mit anderen Teilchen sind entsprechende Prozesse
möglich.

Die Quantenelektrodynamik beschreibt außerdem Vorgänge wie

\begin{itemize}
	\item die Streuung von Elektronen,
	\item die Streuung von Photonen an Elektronen,
	\item Paarbildung,
	\item Annihilation und
	\item Korrekturen des magnetischen Moments.
\end{itemize}

\subsubsection*{Feynman-Diagramme}

Viele Prozesse der Quantenelektrodynamik können mithilfe von
Feynman-Diagrammen%
\index{Feynman-Diagramm}
dargestellt werden.

In solchen Diagrammen werden Teilchen und ihre Wechselwirkungen durch
Linien und Verbindungspunkte symbolisiert.

Typischerweise stehen

\begin{itemize}
	\item gerade Linien für Elektronen oder Positronen,
	\item gewellte Linien für Photonen und
	\item Verbindungspunkte für Wechselwirkungen.
\end{itemize}

Ein Feynman-Diagramm ist jedoch kein genaues Bild einer räumlichen
Teilchenbahn. Es stellt einen mathematischen Beitrag zur
Wahrscheinlichkeit eines Prozesses dar.

\begin{DidacticBox}[Feynman-Diagramme sind keine Fotografien]
	\label{box:feynman_diagramme_keine_fotografien}
	
	Ein Feynman-Diagramm zeigt nicht, wie Teilchen im Raum tatsächlich auf
	kleinen Bahnen fliegen.
	
	Es ist eine grafische Darstellung mathematischer Terme, mit denen die
	Wahrscheinlichkeit eines Wechselwirkungsprozesses berechnet wird.
\end{DidacticBox}

\subsubsection*{Virtuelle Teilchen}

In Berechnungen mit Feynman-Diagrammen treten häufig innere Linien auf.
Sie werden mit sogenannten virtuellen Teilchen%
\index{virtuelles Teilchen} verbunden.

Virtuelle Teilchen sind keine frei beobachtbaren Teilchen, die
vorübergehend direkt gemessen werden könnten.

Sie sind Bestandteile der mathematischen Beschreibung einer
Wechselwirkung. Sie müssen nicht dieselbe Energie-Impuls-Beziehung
erfüllen wie frei beobachtbare Teilchen.

Beobachtbare Elektronen, Positronen und Photonen treten dagegen als
reale Anfangs- oder Endzustände eines Prozesses auf.

\begin{NoteBox}[Vorsicht beim Begriff virtuelles Teilchen]
	\label{box:vorsicht_virtuelle_teilchen}
	
	Die Aussage, ein virtuelles Teilchen leihe sich kurzzeitig Energie aus,
	ist eine verbreitete, aber irreführende Vereinfachung.
	
	Energie bleibt auch in der Quantenfeldtheorie erhalten.
	
	Virtuelle Teilchen sind mathematische Zwischenbeiträge innerhalb einer
	Berechnung und keine direkt beobachtbaren kurzlebigen Teilchen.
\end{NoteBox}

\subsubsection*{Das Vakuum ist nicht einfach leer}

Der Vakuumzustand enthält keine frei beobachtbaren Teilchen. Trotzdem
ist er nicht mit dem klassischen Begriff eines vollkommen leeren
Raumes gleichzusetzen.

Auch im Vakuum sind die Quantenfelder vorhanden. Sie befinden sich in
ihrem niedrigsten Energiezustand.

Die Quanteneigenschaften dieses Zustandes können messbare Auswirkungen
haben. Dazu gehören beispielsweise Korrekturen an Energieniveaus und
am magnetischen Moment des Elektrons.

Solche Effekte werden häufig unter dem Begriff
Vakuumfluktuationen\index{Vakuumfluktuationen} zusammengefasst.

Dabei ist Vorsicht mit anschaulichen Bildern geboten. Das Vakuum sollte
nicht einfach als Raum voller ständig spontan entstehender und wieder
verschwindender klassischer Teilchen verstanden werden.

\begin{DidacticBox}[Quantenfeldtheoretisches Vakuum]
	\label{box:quantenfeldtheoretisches_vakuum}
	
	Das Vakuum ist der Zustand mit der niedrigsten möglichen Energie.
	
	Es enthält keine frei messbaren Teilchen, aber die Quantenfelder selbst
	sind weiterhin vorhanden.
	
	Ihre Quanteneigenschaften können beobachtbare physikalische
	Auswirkungen besitzen.
\end{DidacticBox}

\subsubsection*{Die Dirac-Gleichung bleibt erhalten}

Der Übergang zur Quantenfeldtheorie bedeutet nicht, dass die
Dirac-Gleichung falsch oder überflüssig wird.

Die Gleichung bleibt die grundlegende Bewegungsgleichung des freien
Elektronenfeldes.

Der entscheidende Unterschied besteht darin, dass der Dirac-Spinor nun
als Operatorfeld behandelt wird:

\begin{equation}
	\Psi
	\longrightarrow
	\hat{\Psi}.
\end{equation}

Dadurch können Elektronen und Positronen erzeugt und vernichtet werden.

Die Einteilchen-Dirac-Gleichung ist somit ein wichtiger Grenzfall der
umfassenderen Quantenfeldtheorie.

\begin{PhysicsBox}[Die Dirac-Gleichung in der Feldtheorie]
	\label{box:diracgleichung_in_feldtheorie}
	
	Die Quantenfeldtheorie ersetzt die Dirac-Gleichung nicht.
	
	Sie erweitert ihre Bedeutung:
	
	\begin{itemize}
		\item Aus der Wellenfunktion wird ein quantisiertes Feld.
		\item Elektronen und Positronen werden zu Anregungen dieses Feldes.
		\item Die Anzahl der Teilchen kann sich verändern.
	\end{itemize}
	
	Die mathematische Struktur der Dirac-Gleichung bleibt dabei erhalten.
\end{PhysicsBox}

\subsubsection*{Warum die Feldtheorie notwendig ist}

Die Quantenfeldtheorie löst mehrere Probleme, die in einer
relativistischen Einteilchentheorie auftreten:

\begin{enumerate}
	\item Sie erlaubt eine veränderliche Anzahl von Teilchen.
	\item Sie beschreibt Teilchen und Antiteilchen einheitlich.
	\item Sie behandelt Materie und Strahlung als quantisierte Felder.
	\item Sie ermöglicht die Berechnung von Paarbildung und
	Annihilation.
	\item Sie erklärt Korrekturen, die über die einfache
	Dirac-Gleichung hinausgehen.
\end{enumerate}

Ein besonders bedeutendes Beispiel ist das magnetische Moment des
Elektrons.

Die Dirac-Gleichung liefert zunächst

\begin{equation}
	g
	=
	2.
\end{equation}

Die Quantenelektrodynamik berücksichtigt zusätzlich die Wechselwirkung
des Elektrons mit dem quantisierten elektromagnetischen Feld und
erklärt dadurch die geringe Abweichung vom Wert \(2\).

\subsubsection*{Vom Teilchenbild zum Feldbild}

Der historische Weg zur modernen Beschreibung des Elektrons lässt sich
vereinfacht zusammenfassen:

\begin{equation}
	\begin{aligned}
		\text{klassisches Teilchen}
		&\longrightarrow
		\text{Quantenobjekt}
		\\
		&\longrightarrow
		\text{relativistischer Spinor}
		\\
		&\longrightarrow
		\text{Anregung eines Quantenfeldes}.
	\end{aligned}
	\label{eq:entwicklung_elektronenbild}
\end{equation}

Jede Stufe erweitert die vorherige Beschreibung.

Das klassische Teilchenbild bleibt für viele technische Anwendungen
nützlich. Die Schrödingergleichung beschreibt zahlreiche atomare und
molekulare Vorgänge. Die Dirac-Gleichung ergänzt Relativität, Spin und
Antiteilchen.

Die Quantenfeldtheorie wird notwendig, wenn Teilchen erzeugt,
vernichtet oder durch quantisierte Felder beeinflusst werden.

\begin{DidacticBox}[Mehrere Beschreibungen, kein Widerspruch]
	\label{box:mehrere_beschreibungen_elektron}
	
	Das Elektron kann je nach Fragestellung unterschiedlich beschrieben
	werden:
	
	\begin{itemize}
		\item als klassisches geladenes Teilchen,
		\item durch eine Schrödinger-Wellenfunktion,
		\item durch einen Dirac-Spinor oder
		\item als Anregung eines Quantenfeldes.
	\end{itemize}
	
	Diese Modelle widersprechen sich nicht grundsätzlich. Sie besitzen
	unterschiedliche Gültigkeitsbereiche und unterschiedliche
	Genauigkeiten.
\end{DidacticBox}

\subsubsection*{Bedeutung für die moderne Physik}

Die Verbindung von Quantenmechanik und spezieller Relativitätstheorie
führte nicht lediglich zu einer verbesserten Bewegungsgleichung für
das Elektron.

Sie veränderte das grundlegende Verständnis von Materie.

Teilchen werden in der modernen Physik als quantisierte Anregungen von
Feldern beschrieben. Antiteilchen sind dabei keine mathematischen
Sonderfälle, sondern notwendige Bestandteile relativistischer
Quantenfelder.

Die Quantenelektrodynamik entwickelte sich zu einer der präzisesten
physikalischen Theorien. Sie beschreibt die elektromagnetische
Wechselwirkung von Elektronen, Positronen und Photonen.

\begin{HistoryBox}[Vom Elektron zur Quantenfeldtheorie]
	\label{box:vom_elektron_zur_quantenfeldtheorie}
	
	Die Suche nach einer relativistischen Gleichung für ein einzelnes
	Elektron führte schrittweise zu einer umfassenderen Theorie.
	
	Die negativen Energielösungen wiesen auf Antiteilchen hin. Paarbildung
	und Annihilation zeigten, dass die Teilchenzahl nicht fest sein kann.
	
	Daraus entstand die Vorstellung, dass nicht einzelne Teilchen, sondern
	Quantenfelder die grundlegenden Größen der Theorie sind.
\end{HistoryBox}

\subsubsection*{Zusammenfassung des Kapitels}

Die Schrödingergleichung beschreibt Elektronen zuverlässig im
nichtrelativistischen Bereich. Für hohe Geschwindigkeiten und Energien
musste sie jedoch erweitert werden.

Dirac gelang es, eine Gleichung zu formulieren, die Quantenmechanik und
spezielle Relativitätstheorie miteinander verbindet.

Aus der Dirac-Gleichung folgten

\begin{itemize}
	\item der Spin des Elektrons,
	\item sein magnetisches Moment,
	\item positive und negative Energielösungen sowie
	\item die Existenz des Positrons.
\end{itemize}

Die Möglichkeit von Paarbildung und Annihilation zeigte schließlich,
dass eine feste Teilchenzahl nicht ausreicht.

In der Quantenfeldtheorie werden Elektronen und Positronen daher als
Anregungen eines quantisierten Dirac-Feldes beschrieben. Photonen sind
Anregungen des elektromagnetischen Feldes.

\begin{PhysicsBox}[Vom Elektron zum Quantenfeld]
	\label{box:vom_elektron_zum_quantenfeld}
	
	Die moderne Beschreibung des Elektrons verbindet vier zentrale Ideen:
	
	\begin{enumerate}
		\item Das Elektron besitzt Quanteneigenschaften.
		\item Seine Bewegung muss relativistisch beschrieben werden können.
		\item Es besitzt Spin und ein Antiteilchen.
		\item Es ist eine Anregung eines quantisierten Feldes.
	\end{enumerate}
	
	Damit führt der Weg von der Schrödingergleichung über die
	Dirac-Gleichung unmittelbar zur Quantenfeldtheorie.
\end{PhysicsBox}

\subsubsection*{Fazit}

Die Entwicklung von der Schrödingergleichung über die Dirac-Gleichung
bis zur Quantenfeldtheorie zeigt, wie eng physikalische Prinzipien und
mathematische Strukturen miteinander verbunden sind.

Dirac wollte Quantenmechanik und spezielle Relativitätstheorie in einer
gemeinsamen Gleichung vereinen. Aus den dafür notwendigen
mathematischen Bedingungen folgten jedoch Eigenschaften, die nicht
zusätzlich vorausgesetzt worden waren: der Spin des Elektrons, die
vierkomponentige Struktur seines Zustandes und Lösungen, die auf ein
Antiteilchen hinwiesen.

Das Positron\index{Positron!mathematische Vorhersage} wurde damit
zunächst als mathematische Konsequenz sichtbar und erst später
experimentell nachgewiesen.

\begin{HistoryBox}[Mathematik als Wegweiser]
	\label{box:mathematik_als_wegweiser_dirac}
	
	Die Dirac-Gleichung zeigt, dass Mathematik in der Physik mehr sein kann
	als ein Werkzeug zur Berechnung.
	
	Aus physikalischen Prinzipien und mathematischer Konsequenz entstand
	die Vorhersage einer zuvor unbekannten Naturerscheinung.
\end{HistoryBox}

